% RQ1 Background: Participation Dynamics

\subsection{Participation Challenges in DAOs}

Empirical studies of deployed DAOs consistently document low voter turnout. \citet{fritsch2022analyzing} found that in MakerDAO, fewer than 20 addresses controlled over 50\% of voting power, with typical proposal turnout below 5\%. DeepDAO analytics \citep{deepdao2023analytics} report median turnout rates of 2-8\% across major governance protocols.

This pattern echoes classic collective action problems identified by \citet{olson1965logic}: rational individuals under-invest in public goods when they cannot be excluded from benefits. In DAO governance, the ``public good'' is effective decision-making, and the cost is voter attention.

\subsection{Quorum Design in Practice}

Real-world DAOs employ various quorum mechanisms:

\begin{itemize}
    \item \textbf{Compound}: 4\% of total token supply required for quorum
    \item \textbf{Uniswap}: 4\% quorum with 2.5M UNI proposal threshold
    \item \textbf{Aave}: Tiered quorums based on proposal risk level
    \item \textbf{Optimism}: Token House requires 30\% of votable supply
\end{itemize}

These thresholds were often set through intuition or historical precedent rather than systematic analysis. The consequences of miscalibration are visible: Compound has experienced periods of governance stagnation when participation dropped below quorum, while some DAOs with very low quorums have seen contentious proposals pass with minimal deliberation.

\subsection{Voter Fatigue}

The phenomenon of ``governance fatigue'' has been documented across multiple DAOs \citep{snapshot2021governance}. Contributing factors include:

\begin{itemize}
    \item High proposal frequency overwhelming participant attention
    \item Repetitive or low-stakes votes reducing perceived importance
    \item Lack of meaningful impact from individual votes
    \item Cognitive load of evaluating complex technical proposals
\end{itemize}

Mitigation strategies employed by DAOs include delegation systems (Compound, ENS), conviction voting (Giveth), and governance minimization philosophies (Uniswap).

\subsection{Theoretical Foundations}

Our work builds on several theoretical traditions:

\subsubsection{Rational Voter Models}

The calculus of voting \citep{olson1965logic} models participation decisions as:

\begin{equation}
\text{Vote if } pB - C > 0
\end{equation}

where $p$ is the probability of being pivotal, $B$ is the benefit of preferred outcome, and $C$ is the cost of voting. In large electorates, $p \approx 0$, predicting near-zero turnout---the ``paradox of voting.''

In DAOs, token-weighted voting modifies this calculus: larger holders have higher $p$ and $B$, explaining their disproportionate participation.

\subsubsection{Bounded Rationality}

Real voters do not perform full expected utility calculations. Following \citet{ostrom1990governing}, we model agents with bounded rationality: participation depends on heuristics, habits, and social factors rather than pure calculation.

\subsubsection{Multi-Agent Systems}

Our simulation approach draws from agent-based computational economics \citep{tesfatsion2006handbook} and multi-agent systems \citep{wooldridge2009introduction}. These frameworks enable studying emergent macro-level patterns (aggregate turnout, quorum achievement) from micro-level agent behaviors.

\subsection{Prior Simulation Work}

Limited prior work has specifically addressed DAO participation through simulation. \citet{faqir2021cadcad} used cadCAD for tokenomic modeling but focused on token engineering rather than voting behavior. Our framework addresses this gap by providing dedicated support for participation dynamics modeling.
